\documentclass[11pt]{article}

\usepackage[T1]{fontenc}
\usepackage[utf8]{inputenc}
\usepackage{lmodern}
\usepackage[a4paper,margin=20mm]{geometry}
\usepackage{xcolor}
\usepackage{hyperref}

\definecolor{linkblue}{HTML}{174C77}
\hypersetup{colorlinks=true,urlcolor=linkblue,linkcolor=linkblue,
  pdftitle={Graph Classes and Width Parameters in Lax},pdfauthor={Clemens Kuske}}
\newcommand{\source}[2]{\textsuperscript{\href{https://laxarchive.org/#1/#2.html}{$\nearrow$}}}
\newcommand{\surveychapter}[2]{\setcounter{section}{#1}\section{#2}}
\setlength{\parindent}{0pt}
\setlength{\parskip}{4pt}

\title{Graph Classes and Width Parameters in Lax}
\author{Clemens Kuske}
\date{October 2026}

\begin{document}
\maketitle
\vspace{-1em}

This catalog follows the chapter order of the supplied survey, retaining only
chapters represented by formalized definitions. Chapter numbers match the
survey. Arrows link to source concepts; marked passages open their Lean
encodings alongside the PDF.

\begingroup
\small
\setlength{\parskip}{0pt}
\tableofcontents
\endgroup

\surveychapter{0}{Basic Concepts}

\subsection{Basic graph notions}

These references collect graph classes, their closure properties, graph
parameters, and elementary graph families.

% lax begin Lax199508.GraphClasses
\textbf{Graph classes.} The sparsity lectures represent a class by its finite
simple graphs on canonical vertex sets, giving a common starting point for
the later sparsity notions; see
\source{lax-199508}{Lax199508.GraphClasses}.
% lax end

% lax begin Lax871432.GraphClasses
\textbf{Isomorphism-closed graph classes.} Another treatment makes closure
under isomorphism explicit and provides union and intersection operations;
see \source{lax-871432}{Lax871432.GraphClasses}.
% lax end

% lax begin Lax871432.ClosureProperties
\textbf{Closure properties.} This concept groups closure under induced
subgraphs, minors, contractions, disjoint-union summands, and disjoint union.
In particular, induced-subgraph closure is the usual hereditary property;
see \source{lax-871432}{Lax871432.ClosureProperties}.
% lax end

% lax begin Lax153141.GraphParameters
\textbf{Graph parameters.} This concept gives a shared signature for
natural-number graph parameters and records when two parameters bound one
another up to functions; see
\source{lax-153141}{Lax153141.GraphParameters}.
% lax end

% lax begin Lax68.Paths
\textbf{Paths.} These are the finite path graphs, the simplest connected
examples in the list; every path is a tree. See
\source{lax-68}{Lax68.Paths}.
% lax end

% lax begin Lax68.Stars
\textbf{Stars.} These trees have one centre joined to all other vertices;
their leaves have no edges between them. See
\source{lax-68}{Lax68.Stars}.
% lax end

% lax begin Lax68.Trees
\textbf{Trees.} A tree is connected and acyclic, so both paths and stars
form subclasses of the tree family. See \source{lax-68}{Lax68.Trees}.
% lax end

% lax begin Lax68.Ladders
\textbf{Ladders.} A ladder is a two-row grid, with two paths joined by
corresponding rungs. See \source{lax-68}{Lax68.Ladders}.
% lax end

\subsection{Some special graphs}

The remaining examples follow the survey's special-graph discussion. Where
the classes are nested, triangles lie among maximal outerplanar graphs, and
outerplanar graphs lie among planar graphs. Grids, walls, series-parallel
graphs, and triangulations provide further examples in this chapter.

% lax begin Lax68.Triangles
\textbf{Triangles.} This is the three-vertex complete graph, which is also a
maximal outerplanar graph. See \source{lax-68}{Lax68.Triangles}.
% lax end

% lax begin Lax68.MaximalOuterplanar
\textbf{Maximal outerplanar graphs.} These outerplanar graphs are maximal
under adding edges while preserving outerplanarity; every triangle is an
example. See
\source{lax-68}{Lax68.MaximalOuterplanar}.
% lax end

% lax begin Lax68.Outerplanar
\textbf{Outerplanar graphs.} They admit a crossing-free drawing with every
vertex on the outer face, so this family is contained in the planar graphs. See
\source{lax-68}{Lax68.Outerplanar}.
% lax end

% lax begin Lax68.SeriesParallel
\textbf{Series-parallel graphs.} This class is generated from a single edge
by two-terminal series and parallel constructions. See
\source{lax-68}{Lax68.SeriesParallel}.
% lax end

% lax begin Lax68.Triangulations
\textbf{Triangulations.} These are maximal planar graphs, obtained by
triangulating the faces of a planar embedding. See
\source{lax-68}{Lax68.Triangulations}.
% lax end

% lax begin Lax68.GridsAndWalls
\textbf{Grids and walls.} Rectangular grids and their brick-wall variants
are standard planar examples and recurring structural patterns. See
\source{lax-68}{Lax68.GridsAndWalls}.
% lax end

% lax begin Lax68.Planar
\textbf{Planar graphs.} These graphs admit a drawing in the plane with no
crossing edges; the formalized entry records a straight-line drawing as a
certificate. See \source{lax-68}{Lax68.Planar}.
% lax end

\surveychapter{1}{Perfection, Generalized Perfection, and Related Concepts}

\subsection{Related concepts}

% lax begin Lax9.ChiBoundedness
\textbf{$\chi$-boundedness.} This class property bounds chromatic number as
a function of clique number, a control closely related to graph perfection.
See \source{lax-9}{Lax9.ChiBoundedness}.
% lax end

\surveychapter{2}{Cycles, Chords, and Bridges}

\subsection{Cycle deletion parameters}

These parameters count the deletions needed to remove all cycles, by deleting
vertices or edges.

% lax begin Lax379983.FeedbackVertexNumber
\textbf{Feedback vertex number.} This measures the minimum number of vertices
to remove so the remaining graph is a forest. See
\source{lax-379983}{Lax379983.FeedbackVertexNumber}.
% lax end

% lax begin Lax379983.FeedbackEdgeNumber
\textbf{Feedback edge number.} This is the edge-deletion counterpart: remove
edges until the remaining graph is acyclic. See
\source{lax-379983}{Lax379983.FeedbackEdgeNumber}.
% lax end

\surveychapter{4}{Vertex and Edge Orderings}

\subsection{Orderings and neighborhood measures}

The survey treats elimination and neighborhood-based orderings; these
formalized concepts likewise use linear orders and reachability.

% lax begin Lax199508.Admissibility
\textbf{Admissibility.} Relative to an ordering, this parameter counts short
paths from a vertex to earlier vertices that meet only at their start. See
\source{lax-199508}{Lax199508.Admissibility}.
% lax end

% lax begin Lax199508.ColoringNumbers
\textbf{Generalized coloring numbers.} Weak and strong reachability count
vertices accessible by short paths under a linear order; the parameters
minimize the largest count over all orders. See
\source{lax-199508}{Lax199508.ColoringNumbers}.
% lax end

% lax begin Lax199508.NeighborhoodComplexity
\textbf{Almost-linear neighborhood complexity.} This bounds the number of
distinct neighborhood traces on a vertex set by a near-linear function of
its size. See \source{lax-199508}{Lax199508.NeighborhoodComplexity}.
% lax end

% lax begin Lax9.NeighborhoodComplexity
\textbf{Linear neighborhood complexity.} This companion formulation asks for
a linear bound on distinct neighborhood traces. See
\source{lax-9}{Lax9.NeighborhoodComplexity}.
% lax end

\surveychapter{6}{Forbidden Subgraphs}

\subsection{Forbidden induced subgraphs}

% lax begin Lax214022.Cographs
\textbf{Cographs.} Their forbidden-pattern characterization uses the
four-vertex path; this is the same class referenced in the perfection
chapter. See \source{lax-214022}{Lax214022.Cographs}.
% lax end

% lax begin Lax762056.InducedMinors
\textbf{Induced minors.} This containment variant preserves adjacency and
non-adjacency between retained branch sets. See
\source{lax-762056}{Lax762056.InducedMinors}.
% lax end

% lax begin Lax762056.Grid
\textbf{Excluded induced grid minors.} This class excludes an induced minor
isomorphic to a square grid of the specified size. See
\source{lax-762056}{Lax762056.Grid}.
% lax end

% lax begin Lax710763.GraphClasses
\textbf{Weakly sparse classes.} The formalized property excludes a fixed
complete bipartite graph as a subgraph. See
\source{lax-710763}{Lax710763.GraphClasses}.
% lax end

\subsection{Forbidden minor characterizations}

% lax begin Lax68.GraphMinors
\textbf{Graph minors.} Branch sets package vertex deletion, edge deletion,
and contraction into one containment relation. See
\source{lax-68}{Lax68.GraphMinors}.
% lax end

% lax begin Lax68.GraphTopologicalMinors
\textbf{Topological minors.} Edges of the smaller graph are represented by
internally disjoint paths. See
\source{lax-68}{Lax68.GraphTopologicalMinors}.
% lax end

% lax begin Lax68.GridsAndWalls
\textbf{Grids and walls.} These planar structures also serve as canonical
patterns in minor containment. See \source{lax-68}{Lax68.GridsAndWalls}.
% lax end

% lax begin Lax68.Outerplanar
\textbf{Outerplanar graphs.} The drawing-based class also appears in the
survey's discussion of planar structure. See
\source{lax-68}{Lax68.Outerplanar}.
% lax end

% lax begin Lax68.Planar
\textbf{Planar graphs.} Their forbidden-minor characterization is a central
example in this chapter. See \source{lax-68}{Lax68.Planar}.
% lax end

% lax begin Lax68.Triangulations
\textbf{Triangulations.} These provide maximal planar examples alongside the
planarity and minor references above. See
\source{lax-68}{Lax68.Triangulations}.
% lax end

\surveychapter{8}{Matrices and Polyhedra}

\subsection{Matrix-based graph parameters}

% lax begin Lax153141.MixedMinorNumber
\textbf{Mixed minor number.} This parameter detects fully mixed blocks in
ordered adjacency matrices and is comparable to twin-width. See
\source{lax-153141}{Lax153141.MixedMinorNumber}.
% lax end

\surveychapter{9}{Distance Properties}

\subsection{Distances, shallow minors, and sparse classes}

These concepts use bounded-radius reachability or shallow containment to
describe structural sparsity.

% lax begin Lax199508.UniformQuasiWideness
\textbf{Uniform quasi-wideness.} Large vertex sets contain large
distance-separated subsets after removing a bounded-size separator, uniformly
over the class. See
\source{lax-199508}{Lax199508.UniformQuasiWideness}.
% lax end

% lax begin Lax199508.ShallowMinorDensity
\textbf{Shallow-minor density.} This records edge density after contractions
of bounded depth, for each fixed radius. See
\source{lax-199508}{Lax199508.ShallowMinorDensity}.
% lax end

% lax begin Lax199508.ShallowTopologicalMinors
\textbf{Shallow topological minors.} This bounds the subdivision depth of
paths representing edges in a shallow topological minor. See
\source{lax-199508}{Lax199508.ShallowTopologicalMinors}.
% lax end

% lax begin Lax199508.NowhereDenseClasses
\textbf{Nowhere-dense graph classes.} At each fixed depth, some clique is
excluded as a shallow minor of every graph in the class. See
\source{lax-199508}{Lax199508.NowhereDenseClasses}.
% lax end

% lax begin Lax710763.MonadicDependence
\textbf{Monadically dependent classes.} These classes cannot produce every
graph through the specified first-order transductions with vertex colorings.
See \source{lax-710763}{Lax710763.MonadicDependence}.
% lax end

\clearpage
\surveychapter{10}{Algebraic Compositions and Recursive Definitions}

\subsection{Trees, $k$-trees, and partial $k$-trees}

% lax begin Lax68.Trees
\textbf{Trees.} The recursively describable tree family appears here again
alongside the basic graph reference. See \source{lax-68}{Lax68.Trees}.
% lax end

% lax begin Lax228581.Treewidth
\textbf{Treewidth.} This parameter is defined by tree decompositions and
their largest bag size. See \source{lax-228581}{Lax228581.Treewidth}.
% lax end

\subsection{Series-parallel graphs}

% lax begin Lax68.SeriesParallel
\textbf{Series-parallel graphs.} This class is generated by two-terminal
series and parallel composition. See
\source{lax-68}{Lax68.SeriesParallel}.
% lax end

\subsection{Cographs and width parameters}

% lax begin Lax214022.Cographs
\textbf{Cographs.} The recursive and forbidden-path views of this class are
linked by the survey; its formalized entry also occurs in the earlier
chapters. See \source{lax-214022}{Lax214022.Cographs}.
% lax end

% lax begin Lax271696.CliqueExpr
\textbf{Cliquewidth.} A $k$-expression constructs a graph using $k$ labels;
bounded cliquewidth organizes classes by the number of labels needed. See
\source{lax-271696}{Lax271696.CliqueExpr}.
% lax end

% lax begin Lax9.MergeWidth
\textbf{Merge-width.} This width parameter measures merge sequences using
radius-dependent neighborhoods in partially resolved graphs. See
\source{lax-9}{Lax9.MergeWidth}.
% lax end

% lax begin Lax228581.TwinWidth
\textbf{Twin-width.} This partition-sequence parameter counts red neighbors
created as vertex parts merge, and applies to both sparse and dense classes.
See \source{lax-228581}{Lax228581.TwinWidth}.
% lax end

\surveychapter{11}{Decompositions and Cutsets}

\subsection{Graph decompositions and width}

% lax begin Lax228581.Treewidth
\textbf{Treewidth.} Tree decompositions cover vertices and edges while
keeping each vertex's bags connected; their width is measured by the largest
bag. See \source{lax-228581}{Lax228581.Treewidth}.
% lax end

% lax begin Lax228581.TwinWidth
\textbf{Twin-width.} The sequence of merged vertex parts gives a second
decomposition-based width measure. See
\source{lax-228581}{Lax228581.TwinWidth}.
% lax end

\end{document}
